\documentclass[11pt,reqno]{amsart}

\usepackage[a4paper,margin=1.15in]{geometry}
\usepackage[T1]{fontenc}
\usepackage{lmodern}
\usepackage{amsmath,amssymb,mathtools}
\usepackage{microtype}
\usepackage{enumitem}
\usepackage[hidelinks]{hyperref}
\hypersetup{
  pdftitle={Full Tracial Corners and Finite Central Fibers},
  pdfauthor={}
}

\numberwithin{equation}{section}

\newtheorem{theorem}{Theorem}[section]
\newtheorem{proposition}[theorem]{Proposition}
\newtheorem{lemma}[theorem]{Lemma}
\newtheorem{corollary}[theorem]{Corollary}
\theoremstyle{definition}
\newtheorem{definition}[theorem]{Definition}
\newtheorem{assumption}[theorem]{Assumption}
\theoremstyle{remark}
\newtheorem{remark}[theorem]{Remark}

\newcommand{\cB}{\mathcal B}
\newcommand{\cH}{\mathcal H}
\newcommand{\cM}{\mathcal M}
\newcommand{\cZ}{\mathcal Z}
\newcommand{\id}{\mathbf 1}
\newcommand{\Tr}{\operatorname{Tr}}
\newcommand{\supp}{\operatorname{supp}}
\newcommand{\fin}{\mathrm{fin}}
\newcommand{\cF}{\mathcal F}

\title[Full Tracial Corners]{Full Tracial Corners and Finite Central Fibers\\
\large A mathlib-oriented formalization blueprint}
\author{}
\date{August 2, 2026}

\begin{document}

\begin{abstract}
We isolate the operator-algebraic core of the current stabilizer-region proof.
Let $\cM\subseteq\cB(\cH)$ be a concrete von Neumann algebra and let
$q\in\cM$ be a full projection.  If a complex-linear functional is tracial
on the support corner $q\cM q$ and faithful on the projections under $q$,
then every nonzero central summand of $\cM$ contains a nonzero finite
projection.  Consequently, every member of an explicitly given finite family
of nonzero central projections has a canonical finite witness.  The note
states the exact assumptions and gives the complete algebraic proof that is to
be machine checked.  It deliberately does not claim that the stabilizer/GNS
construction supplies those assumptions, nor does it yet bundle the central
summand as a represented factor to state a machine-checked type-III
conclusion; those are distinct, later formalization layers.
\end{abstract}

\maketitle

\section{Scope, status, and relation to the manuscript}

The current Overleaf manuscript contains the following theorem:
\begin{center}
\texttt{thm:finite\_projections\_exist}
\end{center}
For every nonzero central projection $z$ of a regional von Neumann algebra
$\cM_\Gamma$, the central summand $z\cM_\Gamma$ contains a finite projection.
Its stated route is a full support projection $q$ and a faithful normal
tracial state on $q\cM_\Gamma q$.

This note formalizes only the algebraic implication from those corner
hypotheses to that conclusion.  The status of the statements in this note is
as follows.
\begin{enumerate}[label=\textup{(\roman*)},leftmargin=2.5em]
  \item The definitions and proofs in Sections~\ref{sec:algebraic_data}--
  \ref{sec:finite_fibers} are formalized as the exact abstract Lean target.
  \item The stabilizer/GNS realization listed in
  Section~\ref{sec:stabilizer_boundary} is not claimed here.
  \item A Lean build can certify the abstract target, but it cannot by itself
  certify the manuscript's stabilizer hypotheses.
\end{enumerate}

This separation is essential.  In particular, a finite set of central fibers
does not create the support projection, establish its traciality, or establish
fullness; it merely avoids the additional direct-integral layer needed to
interpret arbitrary measurable fibers.

\section{Algebraic data and conventions}
\label{sec:algebraic_data}

Let $\cH$ be a complex Hilbert space and let
\begin{equation}
  \cM\subseteq\cB(\cH)
  \label{eq:concrete_von_neumann_algebra}
\end{equation}
be a concrete von Neumann algebra.  The intended Lean ambient object is
mathlib's \texttt{VonNeumannAlgebra H}; no new notion of von Neumann algebra
is introduced.

\begin{definition}[Projection and subprojection]
\label{def:projection_subprojection}
A projection in $\cM$ is an operator $e\in\cM$ satisfying
\begin{equation}
  e^*=e,
  \qquad
  e^2=e.
  \label{eq:star_projection}
\end{equation}
For projections $e,f$, write $e\leq f$ when
\begin{equation}
  ef=e,
  \qquad
  fe=e.
  \label{eq:subprojection}
\end{equation}
This two-sided equation is the definition used by the Lean development; it
avoids introducing an order structure before it is needed.
\end{definition}

\begin{definition}[Murray--von Neumann equivalence and finiteness]
\label{def:mvN_finiteness}
For projections $e,f\in\cM$, write $e\sim_{\cM}f$ if there exists
$v\in\cM$ with
\begin{equation}
  v^*v=e,
  \qquad
  vv^*=f.
  \label{eq:mvN_equivalence}
\end{equation}
A projection $e\in\cM$ is \emph{finite} if, for every projection
$f\in\cM$,
\begin{equation}
  e\sim_{\cM}f,
  \qquad
  f\leq e
  \quad\Longrightarrow\quad
  f=e.
  \label{eq:finite_projection}
\end{equation}
\end{definition}

\begin{lemma}[A Murray--von Neumann witness is a partial isometry]
\label{lem:mvN_witness_partial_isometry}
Let $v\in\cB(\cH)$ and let $e$ be a projection.  If $v^*v=e$, then
\begin{equation}
  vv^*v=v.
  \label{eq:mvN_witness_partial_isometry}
\end{equation}
Thus the standard definition \eqref{eq:mvN_equivalence} is equivalent, for
projection endpoints, to the variant which explicitly asks that $v$ be a
partial isometry.
\end{lemma}

\begin{proof}
Put $w=v-ve$.  Since $e=e^*=e^2$ and $v^*v=e$,
\begin{equation}
  w^*w=(1-e)v^*v(1-e)=(1-e)e(1-e)=0.
  \label{eq:partial_isometry_zero_norm}
\end{equation}
The $C^*$-identity in $\cB(\cH)$ implies $w=0$, hence $ve=v$.  Therefore
$vv^*v=v(v^*v)=ve=v$.
\end{proof}

\begin{definition}[Central and full projections]
\label{def:central_full}
A projection $z\in\cM$ is central when
\begin{equation}
  az=za,
  \qquad
  a\in\cM.
  \label{eq:central_projection}
\end{equation}
Thus the central projections are the projections in $\cZ(\cM)$.
A projection $q\in\cM$ is \emph{full} for the present theorem when
\begin{equation}
  z\in\cZ(\cM),\quad z\neq0
  \quad\Longrightarrow\quad
  zq\neq0.
  \label{eq:explicit_fullness}
\end{equation}
For a von Neumann algebra, this is equivalent to the usual condition that the
central support of $q$ is $\id$; the explicit formulation
\eqref{eq:explicit_fullness} is chosen because it is precisely what the proof
uses.  Compare Blackadar, Definition III.1.1.5, as cited in the current
manuscript's fullness discussion.
\end{definition}

\begin{definition}[Semifiniteness in projection form]
\label{def:semifinite_projection_form}
The algebra $\cM$ is called semifinite in the projection form used here if
\begin{equation}
  \forall z\in\cZ(\cM),\quad z\neq0
  \quad\Longrightarrow\quad
  \exists e\in\cM:
  \ e\neq0,\ e\leq z,\ e\text{ is finite}.
  \label{eq:semifinite_projection_form}
\end{equation}
This is the exact conclusion of the main theorem under formalization.  It is
the standard projection-theoretic obstruction to a type-III summand; see the
type definitions in Takesaki~I, Definition V.1.17, and Blackadar,
Section III.1.4.
\end{definition}

\section{The tracial-corner hypothesis}
\label{sec:tracial_corner}

The manuscript obtains its functional by compressing a GNS vector state.  The
abstract algebraic proof needs less than normality and less than faithfulness
on all positive elements.  The following finite-projection test is sufficient
and is a formal consequence of a faithful positive tracial state on $q\cM q$.

\begin{assumption}[Faithful tracial projection functional]
\label{ass:faithful_tracial_projection_functional}
Fix a projection $q\in\cM$.  There is a complex-linear functional
\begin{equation}
  \tau:\cB(\cH)\longrightarrow\mathbb C
  \label{eq:trace_functional}
\end{equation}
whose restriction to the corner $q\cM q$ has the following properties.  We
represent $a\in q\cM q$ in the ambient algebra by the three conditions
\begin{equation}
  a\in\cM,
  \qquad
  qa=a=aq.
  \label{eq:corner_element}
\end{equation}
Then:
\begin{enumerate}[label=\textup{(T\arabic*)},leftmargin=3em]
  \item for all $a,b\in\cM$ satisfying \eqref{eq:corner_element},
  \begin{equation}
    \tau(ab)=\tau(ba);
    \label{eq:trace_identity}
  \end{equation}
  \item for every projection $r\in\cM$ with $r\leq q$,
  \begin{equation}
    \tau(r)=0
    \quad\Longrightarrow\quad
    r=0.
    \label{eq:faithful_on_corner_projections}
  \end{equation}
\end{enumerate}
\end{assumption}

\begin{remark}[Why this is faithful to the manuscript]
\label{rem:faithful_state_implies_assumption}
If $\tau_q$ is a faithful positive tracial state on $q\cM q$, define
$\tau(x)=\tau_q(qxq)$ on the operators relevant below.  Any projection
$r\leq q$ lies in $q\cM q$, and positivity plus faithfulness of $\tau_q$
gives \eqref{eq:faithful_on_corner_projections}.  The trace identity
\eqref{eq:trace_identity} holds when both arguments satisfy
\eqref{eq:corner_element}, because then their products are computed in
$q\cM q$.  No tracial identity on arbitrary elements of $\cM$ is asserted or
needed.  Thus
Assumption~\ref{ass:faithful_tracial_projection_functional} is not a stronger
replacement for the manuscript's corner state; it is a deliberately minimal
interface for the finite-projection argument.
\end{remark}

\section{Finiteness below a tracial support}
\label{sec:corner_finiteness}

\begin{lemma}[Difference of nested projections]
\label{lem:projection_difference}
Let $f,e,q\in\cM$ be projections with $f\leq e\leq q$.  Then $e-f$ is a
projection in $\cM$ and
\begin{equation}
  e-f\leq q.
  \label{eq:difference_below_q}
\end{equation}
\end{lemma}

\begin{proof}
The relation $f\leq e$ gives $fe=f=ef$.  Hence
\begin{align}
  (e-f)^2
  &=e^2-ef-fe+f^2 \\
  &=e-f,
  \label{eq:difference_idempotent}
\end{align}
and $(e-f)^*=e-f$.  Since $e,f\in\cM$ and $\cM$ is a $*$-subalgebra,
$e-f\in\cM$.  Transitivity of the two relations in
\eqref{eq:subprojection} gives $fq=f=qf$, while $eq=e=qe$.  Therefore
\begin{align}
  (e-f)q&=eq-fq=e-f,\\
  q(e-f)&=qe-qf=e-f,
  \label{eq:difference_subprojection}
\end{align}
which is \eqref{eq:difference_below_q}.
\end{proof}

\begin{lemma}[Murray--von Neumann witnesses lie in the corner]
\label{lem:mvN_witness_in_corner}
Let $f,e,q\in\cM$ be projections with $f\leq e\leq q$.  If $v\in\cM$
satisfies
\begin{equation}
  v^*v=e,
  \qquad
  vv^*=f,
  \label{eq:corner_mvN_witness}
\end{equation}
then
\begin{equation}
  qv=v=vq,
  \qquad
  qv^*=v^*=v^*q.
  \label{eq:corner_supported_witness}
\end{equation}
In particular, both $v$ and $v^*$ are elements of $q\cM q$ in the sense of
\eqref{eq:corner_element}.
\end{lemma}

\begin{proof}
By Lemma~\ref{lem:mvN_witness_partial_isometry},
$vv^*v=v$.  Hence
\begin{align}
  ve&=v(v^*v)=vv^*v=v,\\
  fv&=(vv^*)v=v.
  \label{eq:initial_final_support_witness}
\end{align}
Since $f\leq e\leq q$, transitivity gives $f\leq q$.  Therefore
\begin{align}
  qv&=q(fv)=(qf)v=fv=v,\\
  vq&=(ve)q=v(eq)=ve=v.
  \label{eq:corner_support_witness}
\end{align}
Taking adjoints gives the two identities for $v^*$ in
\eqref{eq:corner_supported_witness}.  Membership in $\cM$ follows from
$v\in\cM$ and $*$-closure.
\end{proof}

\begin{lemma}[Finite subprojections of a faithful tracial corner]
\label{lem:finite_subprojection_of_tracial_corner}
Under Assumption~\ref{ass:faithful_tracial_projection_functional}, every
projection $e\in\cM$ with $e\leq q$ is finite in $\cM$.
\end{lemma}

\begin{proof}
Let $f\in\cM$ be a projection with $f\leq e$ and
$e\sim_{\cM}f$.  Choose $v\in\cM$ satisfying
  \begin{equation}
    v^*v=e,
    \qquad
  vv^*=f.
  \label{eq:corner_equivalence_witness}
  \end{equation}
Lemma~\ref{lem:mvN_witness_in_corner} puts both $v^*$ and $v$ in $q\cM q$.
By Lemma~\ref{lem:projection_difference}, $e-f$ is a projection below $q$.
Linearity and the corner traciality axiom (T1), applied to $v^*$ and $v$, give
\begin{equation}
  \begin{aligned}
    \tau(e-f)
    &=\tau(e)-\tau(f)\\
    &=\tau(v^*v)-\tau(vv^*)\\
    &=0.
  \end{aligned}
  \label{eq:trace_difference_zero}
\end{equation}
Faithfulness on the projections below $q$ gives $e-f=0$, hence $e=f$.  This
is exactly the finiteness criterion \eqref{eq:finite_projection}.
\end{proof}

\section{Fullness produces finite central summands}
\label{sec:semifiniteness}

\begin{lemma}[A central slice of the corner]
\label{lem:central_slice}
Let $z\in\cZ(\cM)$ be a projection and let $q\in\cM$ be a projection.
Then $zq$ is a projection in $\cM$ and
\begin{equation}
  zq\leq q,
  \qquad
  zq\leq z.
  \label{eq:central_slice_subprojection}
\end{equation}
\end{lemma}

\begin{proof}
Centrality gives $zq=qz$.  Consequently,
\begin{align}
  (zq)^*&=q^*z^*=qz=zq,\\
  (zq)^2&=zqzq=z^2q^2=zq.
  \label{eq:central_slice_projection}
\end{align}
The two relations in \eqref{eq:central_slice_subprojection} follow by
multiplying $zq$ on either side by $q$ and by $z$, respectively.
\end{proof}

\begin{theorem}[Full tracial corner implies projection-form semifiniteness]
\label{thm:full_tracial_corner_semifinite}
Let $\cM\subseteq\cB(\cH)$ be a concrete von Neumann algebra.  Let
$q\in\cM$ be a full projection, and assume
Assumption~\ref{ass:faithful_tracial_projection_functional}.  Then $\cM$ is
semifinite in the sense of Definition~\ref{def:semifinite_projection_form}.
\end{theorem}

\begin{proof}
Let $z\in\cZ(\cM)$ be a nonzero projection.  Fullness gives
\begin{equation}
  zq\neq0.
  \label{eq:central_slice_nonzero}
\end{equation}
By Lemma~\ref{lem:central_slice}, $zq$ is a projection in $\cM$ below $q$
and below $z$.  Lemma~\ref{lem:finite_subprojection_of_tracial_corner}
therefore says that $zq$ is finite.  It is the nonzero finite projection
required by \eqref{eq:semifinite_projection_form}.
\end{proof}

\section{Finite central fibers}
\label{sec:finite_fibers}

\begin{assumption}[Finite central-fiber presentation]
\label{ass:finite_central_fibers}
Let $I$ be a finite set and let $(z_i)_{i\in I}$ be nonzero central
projections in $\cM$ such that
\begin{equation}
  z_i z_j=0\quad (i\neq j),
  \qquad
  \sum_{i\in I}z_i=\id.
  \label{eq:finite_central_partition}
\end{equation}
The factor condition on $z_i\cM$ is not needed to produce finite projections;
it is needed only to use the customary phrase ``fiber factor.''
\end{assumption}

\begin{corollary}[Finite central fibers carry finite projections]
\label{cor:finite_fiber_finite_projection}
Under the hypotheses of
Theorem~\ref{thm:full_tracial_corner_semifinite} and
Assumption~\ref{ass:finite_central_fibers}, each central fiber $z_i\cM$
contains the nonzero finite projection $z_iq$.
\end{corollary}

\begin{proof}
Apply Theorem~\ref{thm:full_tracial_corner_semifinite} to the nonzero central
projection $z_i$.  Its proof identifies the finite witness as $z_iq$, which
lies under $z_i$ by Lemma~\ref{lem:central_slice}.
\end{proof}

\begin{remark}[Type-III conclusion and the corner API]
\label{rem:typeIII_corner_api}
For a factor represented on the Hilbert space $z_i\cH$, a nonzero finite
projection rules out type III by the definition of type III.  A Lean statement
with that exact wording needs a bundled construction of the corner
$z_i\cM$ acting on $z_i\cH$, including preservation/reflection of
Murray--von Neumann finiteness.  The present package records the stronger
ambient statement \eqref{eq:semifinite_projection_form} first.  The finite
fiber corollary will be upgraded to an explicit \texttt{not\_typeIII} theorem
only after that corner API is checked.  This is a scope boundary, not a
mathematical weakening.
\end{remark}

\section{Boundary with the stabilizer/GNS argument}
\label{sec:stabilizer_boundary}

To apply Theorem~\ref{thm:full_tracial_corner_semifinite} to the current
manuscript, a later formalization must construct data satisfying the following
four obligations.
\begin{enumerate}[label=\textup{(S\arabic*)},leftmargin=3em]
  \item From a decreasing exhaustion of local support projections, construct
  a strong-limit projection $q\in\cM_\Gamma$.
  \item Prove that the compressed vector functional on
  $q\cM_\Gamma q$ is tracial.  The manuscript uses finite-dimensional matrix
  trace cyclicity, bounded strong continuity of multiplication, and
  Kaplansky density.
  \item Prove faithfulness on the corner using the dense GNS orbit generated
  by the region and its complement.
  \item Prove fullness by showing that a central projection $z$ with $zq=0$
  annihilates the same dense GNS orbit and hence vanishes.
\end{enumerate}

The source locations in the current manuscript are the support-limit argument
near Lemma~\texttt{lem:q\_exists}, traciality in
Proposition~\texttt{prop:normal-trace}, faithfulness in
Proposition~\texttt{prop:faithful}, fullness in
Proposition~\texttt{prop:fullness}, and the final conclusion in
Theorem~\texttt{thm:finite\_projections\_exist}.  The first two steps require
operator-topological machinery beyond the abstract package in this
investigation.

\section{Lean API and source map}
\label{sec:api_source_map}

The Lean package at
\path{formalization_finite_fibers/} exports the following public declarations:
\begin{itemize}[label={},leftmargin=0pt,itemsep=0pt,topsep=.25em]
  \item \texttt{FiniteFiber.}\allowbreak\texttt{IsProjectionIn}
  \item \texttt{FiniteFiber.}\allowbreak\texttt{IsSubprojection}
  \item \texttt{FiniteFiber.}\allowbreak\texttt{IsMvNEquivalent}
  \item \texttt{FiniteFiber.}\allowbreak\texttt{IsFiniteProjection}
  \item \texttt{FiniteFiber.}\allowbreak\texttt{IsCornerElement}
  \item \texttt{FiniteFiber.}\allowbreak\texttt{IsPartialIsometry}
  \item \texttt{FiniteFiber.}\allowbreak\texttt{FaithfulTracialCorner}
  \item \texttt{FiniteFiber.}\allowbreak\texttt{mvN\_witness\_in\_corner}
  \item \texttt{FiniteFiber.}\allowbreak\texttt{central\_mul\_finite\_of\_faithful\_tracial\_corner}
  \item \texttt{FiniteFiber.}\allowbreak\texttt{finite\_central\_fiber\_canonical\_witness}
  \item \texttt{FiniteFiber.}\allowbreak\texttt{full\_tracial\_corner\_semifinite}
\end{itemize}
Its upstream mathlib bases are:
\begin{itemize}[label={},leftmargin=0pt,itemsep=0pt,topsep=.25em]
  \item \texttt{Mathlib.Analysis.}\allowbreak\texttt{VonNeumannAlgebra.Basic}
  \item \texttt{Mathlib.Algebra.Star.}\allowbreak\texttt{StarProjection}
\end{itemize}

The load-bearing mathematical sources are:
\begin{enumerate}[label=\textup{[\arabic*]},leftmargin=2.5em]
  \item the current local manuscript, Appendix theorem
  \texttt{thm:finite\_projections\_exist} and its immediately preceding
  support-corner proof;
  \item B.~Blackadar, \emph{Operator Algebras}, Definition III.1.1.5 for
  central support, and the finite/type-III projection definitions cited by
  the manuscript in Section III.1.4;
  \item M.~Takesaki, \emph{Theory of Operator Algebras I}, Definition V.1.17
  for the projection-theoretic type classification vocabulary.
\end{enumerate}

\begin{remark}[Verification status]
\label{rem:verification_status}
This TeX note is an informal blueprint.  Its compilation checks typesetting
and references only.  The abstract Lean package has been checked with
\texttt{lake build FiniteFiber}; the accompanying
\texttt{AxiomAudit.lean} in its \texttt{tests/} directory records the public
theorem axiom audit.  That check establishes only the named abstract
support-corner theorem.  The stabilizer/GNS application remains separately
auditable.
\end{remark}

\end{document}
